\subsection{Human Dignity, and One Indicator That Can Be Audited}
\label{sec:5.1.3}
Respecting human dignity through technology means applying AI in ways that strengthen human capability and community ties, in a domain where the same tools can as easily substitute for them. The applications are easy to name and mostly easy to believe: a system that surfaces wage gaps and retaliatory scheduling across large volumes of employment data turns an information asymmetry that currently favors employers into something workers can act on. The same predictive capacity, aimed at the material conditions that let authoritarian movements gain ground, has a poor track record — forecasting political instability from social and economic indicators is an old field with modest results, and this book does not propose building on it.

Naming an indicator is harder than promising one. The one this section names is narrow in what it measures and not at all narrow in what it means, because what it measures is a feature of the definition itself.

\runin{Measure targets per hour against reviewers available}
Wherever an AI system generates candidate decisions that a human is nominally required to approve — welfare eligibility, deportation referrals, content removals, military targets — the ratio between the rate at which the system produces candidates and the number of qualified people available to examine them is a measurable, auditable property of the deployment. The ratio requires no access to the model and no judgment about intent. It says directly whether the human check is a check or a formality, because past a certain ratio the reviewer's only available behavior is approval. Section~\ref{sec:6.4.1} shows what that ratio looks like when it has been allowed to run: a reviewer with twenty seconds per case, describing themselves as a stamp.

What the ratio is not is a general measure of responsible automation, and reading it as one inverts it. Most of what automation exists for is work nobody wants a person to see case by case. An order system that stopped ten thousand times to ask would have reproduced the job rather than done it, and a coding assistant requiring approval for every edit has moved the work instead of removing it. Attention is scarce, and spreading it across thousands of confirmations is how it stops being available on the occasion that needed it. For those deployments the right number of seconds per item is zero, and the indicator has nothing to say about them.

What it is, and the scope has to be stated narrowly or the rest of this section will be read for more than it says, is an instrument for testing the truth of a claim. It bites only where individualized human judgment is what makes the act permissible — where a law, a policy, or a public justification says that a person decided this case. There the ratio says whether that person could have. Where no such claim is made it measures nothing. Where the claim is made and the ratio is bad, what has been established is that the claim is false, and that is the whole of what has been established. An institution whose ratio fails has lost an alibi rather than been stopped, and every act it was performing remains available to it by any route that does not run through the claim.

One consequence of taking the claim seriously is usually read past. Where review cannot keep pace, what ought to give is the rate of action rather than the depth of the review. Calling deliberation a bottleneck treats the output rate as the fixed quantity and moral attention as the inefficiency, when for an irreversible act the capacity to deliberate is the speed limit and the tempo is what has to come down to meet it. Generating candidates faster does not manufacture permission to act on them faster. That is an argument and not a mechanism: an institution that accepts it and slows down has been persuaded, and the number did not do it.

The welfare case is not hypothetical and the record on it is older than the current systems. Virginia Eubanks's study of automated eligibility in Indiana documents what the ratio produces when it is allowed to run: a system that converted any error anywhere in an application into a denial for “failure to cooperate,” refused more than a million applications in three years, a 54 percent increase over the three years before it, and was ended by contract termination after the damage was done \autocite{eubanks2018automating}. The denial reason was a category that could not be contested, because it named no specific failure to answer.

That case also shows the only mode of operation this instrument has been demonstrated to have, which is publication. Nothing about the Indiana system was secret. What was missing was any party positioned to compare the rate of decisions against the capacity to review them, and to say so where it would be read. The ratio is computed, published, and then acted on by a legislature, a court, a union, a caseworkers' association or a reporter, or not acted on at all. It binds nobody. A version that bound somebody would be a regulation, and this is not one.

Across those occasions the indicator generalizes, because the failure it detects is architectural rather than political. It applies to a benefits agency and a targeting cell on identical terms, and an institution that refuses to publish it is telling you something.

What the ratio measures is the fourth feature of the definition in section~\ref{sec:2.1.2}. Acceleration past the rate at which anything can be checked is what opens the gap between an institution's account of the world and the world, and this ratio is that acceleration with a number on it. So the thirteen thousand targets in thirty-eight days in section~\ref{sec:6.4.1} are not an aside about a review step that turned out to be thin. They are the definition being satisfied, and satisfied in public: the figure is not a leak, it is an achievement, given out by the office whose achievement it was. A rate published as a triumph is the second feature arriving alongside the fourth, which is what the two look like when they are working together.

\runin{What it does not reach}
The list of what the ratio does not reach is longer than the list of what it does, and in the military case the gap is widest. It says nothing about area effects, where the question of who was individually selected does not arise. It says nothing about ordnance choice, so a deployment that satisfies the ratio and then strikes an approved target with a weapon that takes the building has passed. It says nothing about the criterion. A classifier scoring people by estimated affiliation and understood inside the unit to be wrong about one time in ten has a defect that review time does not touch, because the reviewer is checking whether the machine applied the criterion and not whether the criterion should be applied to anybody: section~\ref{sec:6.4.1}'s error rate is not a review-time problem, and an indicator that improved while the criterion stayed would be measuring the wrong failure. And it says nothing to a campaign that drops the individualized-judgment claim rather than satisfying it. The ratio tests a stated claim, so an institution that stops stating it has stepped outside the instrument, at a political price that varies with who is watching and is sometimes nothing.

The last omission is the one that should worry a reader who accepted section~\ref{sec:3.6}. A constraint on a rate survives the decomposition objection there because a rate is a fact about the assembly rather than about any request in it. This ratio is that constraint's auditable proxy, and the proxy is measured at the review step. A pipeline whose model does enrichment, collation and ranking upstream of any decision a human is nominally required to approve yields no ratio to compute, and the assembly goes on doing what the assembly does. The instrument is weaker than the floor item it stands in for, and the difference between them is exactly the first interception position.

Iraq makes the limit concrete, and it runs the opposite way from the case this section has been using. Human Rights Watch's review of the 2003 war found that fifty strikes aimed at named members of the Iraqi leadership killed none of the intended targets and killed dozens of civilians, while cluster munitions, overwhelmingly the ground-launched ones fired into populated areas, killed or wounded more than a thousand \autocite{hrw2003offtarget}. The target list was the part that failed at its own job. The mass lethality came from a weapon that requires no list, no candidate decision and no reviewer, so there is nothing in it for this indicator to divide. Mass killing has never needed individualized targeting, and the instrument proposed here is defined only over the part that claims to use it. From which a harder possibility follows: a rate constraint may displace rather than reduce. A force that cannot approve targets individually fast enough retains the option of means that do not require approving them individually, and on that path the ratio improves while the toll does not. Anybody proposing this indicator owes that possibility out loud, and this paragraph is it.

\runin{The second feature, run on this proposal}
Section~\ref{sec:2.1.2}'s second feature applies to what this section is recommending, and it is better run here than left for somebody else. Aestheticization of the metric is the measure substituting for the thing measured, with the substitution experienced as rigor. Publish this ratio as an audit standard and a deployment that clears it acquires a certificate of meaningful human control — a number, defensible, checkable, and about the wrong thing, since what it certifies is that reviewers had time and not that anything was reviewed or that the criterion was worth applying. The unmeasured remainder is the two paragraphs above, and the remainder is where the deaths are.

Simon Zhuang and Dylan Hadfield-Menell give the formal version, and it tells against this proposal rather than for it: optimizing a proxy defined over a subset of the attributes a principal actually values drives overall utility arbitrarily low once resources are constrained, which is a result about proxies of this exact shape rather than a general caution about metrics \autocite{zhuang2020consequences}. This indicator is a proxy defined over a strict subset. The defense available to it is narrow and it is real: it is a proxy for the truth of a claim rather than for a good, and being wrong about whether an institution's account of itself holds up is a smaller error than being wrong about what a life is worth. It survives on that distinction and on one practice. The ratio has to be reported next to what it excludes, and an institution that publishes the number without the exclusions has produced the second feature rather than a check on the fourth.

What the indicator is for is a person auditing one institution, and what it changes is the odds, by making one false claim about a process checkable early enough that a response is still cheap. It is not the first component of something larger. A system that recognizes authoritarian tactics as such is not a capability this book proposes, and the indicator does not become one by being repeated.

None of this works if the system stays a tool wielded by whoever already holds power over workers, welfare recipients, or vulnerable populations. Shifting leverage toward the people affected has to be built in as a design commitment; left to good intentions, it is a side effect.
